% =============================================================================
% §1 Introduction  — Rapid Review lives or dies here (first two pages).
% Target: 1.8--2.0 pages including Fig.~1 and the contribution list.
% 4.5/5 bar: problem is ASPLOS (OS residency + memory/storage hierarchy),
% novelty is scoped, contributions are falsifiable, first two pages self-contained.
% 填写: 把 \num{...} 换成终表数字；替换 Fig.~1 框图；不要改贡献句的范围。
% =============================================================================
\section{Introduction}
\label{sec:intro}

Billion-scale approximate nearest neighbor search (ANNS) must scale two
\emph{independent} dimensions: \emph{corpus capacity} and \emph{query
throughput}. Adding query hosts does not, by itself, shrink the full-precision
graph and embedding store; growing that store does not, by itself, create more
CPU and I/O. Production serving therefore faces a deployment question that
neither ``use a better graph'' nor ``buy a faster SSD'' answers:

\begin{quote}
\emph{When query load grows, can we add compute without replicating the cold
corpus, and without lifting that corpus into DRAM-priced memory?}
\end{quote}

Two published deployment points dominate the literature, and each taxes one
of those dimensions. Block-SSD graph systems such as
DiskANN~\cite{diskann,freshdiskann} and its pipelined descendants~\cite{pipeann}
are excellent at hiding NAND latency behind asynchronous block I/O, product
quantization (PQ) navigation, and a per-host node cache. Their common serving
unit, however, is still \emph{CPU $+$ DRAM working set $+$ a local SSD index}.
Scaling throughput by $R$ query workers therefore typically
replicates provisioned flash ($\propto R\cdot D$) or shards the graph and pays
cross-shard fan-out. Remote block access (NVMe-oF) moves the cable, not the
tax: each host still maintains its own page/buffer cache. Distributed
serving~\cite{distributedann} shows that a single logical DiskANN graph can
live behind a shared KV store, but the interface remains block/KV, not a
unified load/store address space.

The complementary point is full-in-memory CXL-DRAM ANNS.
CXL-ANNS~\cite{cxlanns,cxlanns-tocs} and Cosmos~\cite{cosmos} put the graph
and embeddings in a CXL memory pool so that search never stalls on flash.
That design eliminates NAND misses, and a pooled HDM \emph{can} be one copy
--- but it prices the \emph{cold tail} at DRAM. For a corpus of size $D$ whose
hot working set $H$ satisfies $H\ll D$, the stranded-capacity problem is
structural~\cite{bauhaus,secondtier}: most vectors are rarely touched and
still occupy a DRAM-class device.

Flash-backed CXL memory (CXL-SSD) is the missing third point
(Figure~\ref{fig:intro-position}). A Type-3 memory-semantic SSD exposes NAND
capacity and a small device/host DRAM window through \texttt{load}/\texttt{store}
on a CXL.mem address space~\cite{skybyte,fromblocktobyte,cmm-h,cxl30}. Relative
to replicated DiskANN, the opportunity is to stop copying $D$ with every added
query host; relative to CXL-ANNS/Cosmos, it is to price the cold corpus at
flash. Under CXL 3.x Shared FAM/GFAM, multiple hosts can map the same region
concurrently~\cite{cxl30} --- a \emph{deployment} path to one-copy serving, not
a mechanism we claim to have evaluated. We state the three CXL conditions
explicitly so they are not conflated: (i)~direct-attach CXL-SSD already
provides flash-priced memory semantics on one host; (ii)~CXL 2.0 pooling
shares a device, not a concurrently mapped HDM region; (iii)~only Shared FAM
is a standard premise for multi-host one-copy.

\begin{figure*}[t]
  \centering
  \phbox{0.98\textwidth}{3.6cm}{%
    Fig.~1 placeholder --- draw three columns:\\
    (a)~replicated local-SSD DiskANN ($R\times D$ flash, per-host cache);\\
    (b)~full-in-memory CXL-DRAM ANNS ($D$ at DRAM price);\\
    (c)~\sys{}: graph/PQ in host DRAM, hot records in a bounded CXL-DRAM
    window, full-precision corpus on CXL-SSD, optional Shared-FAM caption.}
  \caption{Three deployment points for billion-scale graph ANNS, and the
    \sys{} runtime that sits on the third. CXL-SSD is an \emph{opportunity}
    (flash-priced corpus, memory-semantic hot path, one-copy under Shared FAM).
    The \emph{system} is residency and stall control on that tier.
    \textbf{Fill-in:} replace this box with a 3-column deployment drawing plus
    a compact runtime overlay on (c). Do not draw a 160\,GB/s device-internal
    bus as a measured link.}
  \Description{Three-column comparison of replicated SSD ANNS, full CXL-DRAM
    ANNS, and the proposed flash-backed CXL-SSD system.}
  \label{fig:intro-position}
\end{figure*}

CXL-SSD does \emph{not} make NAND as fast as DRAM. A cold vector load on this
tier costs tens of microseconds --- two orders of magnitude above a CXL-DRAM
or host-DRAM hit (Figure~\ref{fig:mot-cliff}). Treating the device as slightly
slow memory, or as ``mmap DiskANN'', fails in four systematic ways that we
measure in~\S\ref{sec:motivation}: (1)~the miss cliff dominates end-to-end
search once any material fraction of loads miss; (2)~OS-style page admission
is a poor match for high-dimensional graph adjacency and inflates NAND reads;
(3)~synchronous or high-coverage graph prefetch pollutes the small DRAM window
and can \emph{reduce} QPS; (4)~a blocking load/store miss collapses device
queue depth to $\approx 1$, so a scheduler designed for \texttt{io\_uring}
depth can become a pessimization until fills are made asynchronous. Widening
search (\texttt{efSearch} / beam) to ``buy recall'' further cools the working
set and couples quality to stall.

This paper presents \sys{}, a graph-ANNS \emph{runtime} whose service
goal is to make CXL-SSD \emph{invisible on the search critical path}.
The host should score as if the working set already lived in a small
CXL-DRAM window; the corpus stays on flash because that window cannot
hold $D$. Prefetch, page admission, bandwidth caps, and continuous
batching exist to finish NAND fills \emph{before} a hop scores---not to
make blocking \texttt{load}/\texttt{store} slightly faster.
\sys{} does not replace DiskANN's block-SSD path and does not assume
the corpus fits in CXL-DRAM. A single walk cannot hide hop $i{+}1$
until hop $i$ scores (pointer-chasing); inter-query continuous issue is
what keeps device depth up while one walk is stuck.
The runtime has five coupled modules
(\S\ref{sec:design}): semantic placement; bounded window admission;
entry / soft-pin (not a thrashing LFU hot set); stall-aware search
control; and continuous fetch-level batching over asynchronous
memory-semantic reads.

To our knowledge, this is the first \emph{graph-ANNS serving system} that
directly operates over an SSD-backed CXL address space and jointly manages
ANNS-specific residency, promotion, and search stalls. We do \emph{not} claim
to be the first work to attach a CXL-SSD to a vector database, nor that we
have deployed Shared FAM in production.

This paper makes four contributions:

\begin{enumerate}
  \item \textbf{A third deployment point, stated with its limits.} We
    separate the CXL-SSD \emph{opportunity} (flash-priced $D$, a small
    memory-semantic window, optional one-copy) from the \emph{system}
    (hide NAND from the ANNS critical path), and we name
    the CXL 2.0 vs.\ 3.x conditions under which one-copy is even
    well-formed~(\S\ref{sec:intro}, \S\ref{sec:background}).
  \item \textbf{A measured pathology of graph ANNS on this tier.} We show
    that naive page caches, blind prefetch, unconstrained search width, and
    blocking load/store each amplify stall, and we turn the negative results
    into design constraints~(\S\ref{sec:motivation}).
  \item \textbf{A runtime that tries to hide flash from graph search.}
    \sys{} places graph/PQ off the window, admits few pages, prefetches
    the next hop off the scoring path, and batches fills across queries
    under an iso-recall $L$~(\S\ref{sec:design}, \S\ref{sec:impl}).
  \item \textbf{An iso-recall evaluation of hide, not just QPS.}
    At matched recall we report critical-path wait, window-hit on scored
    vectors, QPS vs.\ demand and vs.\ an in-window oracle, and a
    same-machine DiskANN/PipeANN residual~(\S\ref{sec:eval}).
\end{enumerate}

\paragraph{Scope we do not claim.}
The prototype is an FPGA Type-2/BAR endpoint (\texttt{mem2nvme}/\texttt{nvmex})
plus a software memory-semantic node (\texttt{/dev/vmem0}, \texttt{vmem\_sw})
in front of a 1.92\,TB NVMe. Device HPS and CPU writes into the 32\,GiB BAR
are not live; we do not publish BAR-path numbers. A Montage-style CXL-DRAM
DAX window is \emph{absent} on the measured boot---the 1\,GiB search window
is host DRAM (\texttt{mbind}). We do not report 160\,GB/s device-side copy,
product CXL.mem SSD, or Shared-FAM one-copy. Cost claims are qualitative.

The rest of the paper is organized as follows.
\S\ref{sec:background} reviews graph ANNS and CXL memory-semantic flash.
\S\ref{sec:motivation} presents the findings that constrain the design.
\S\ref{sec:design}--\S\ref{sec:impl} describe \sys{}.
\S\ref{sec:eval} evaluates it.
\S\ref{sec:related} and~\S\ref{sec:discuss} place the work and state limits.
